\section{\textbf{Knowledge Fusion as a Pathway to Knowledge Invention}}
\label{sec:knowledge_fusion}

Knowledge is distributed across disciplines, institutions, datasets,
terminologies, and historical periods. Two scientific communities may study
related structures without sharing concepts, cite separate literatures while
describing interacting mechanisms, or possess results that become
consequential only when interpreted together. This fragmentation creates a
space in which the components of a new insight may already be publicly
available even though the insight itself is absent from collective knowledge.
Knowledge fusion is the process by which an intelligent system identifies and
constructs such a missing relation.

The possibility of deriving new hypotheses from disconnected literatures has
a substantial history. Swanson described \emph{undiscovered public knowledge}
through the example of separately published biomedical findings whose latent
connection suggested a previously unrecognized hypothesis
\cite{swanson1986fish}. Later work on literature-based discovery generalized
this idea into computational methods for identifying implicit relations across
scientific corpora. Studies of scientific innovation likewise show that
atypical combinations of prior work are associated with unusually high impact
\cite{uzzi2013atypical}, while research that bridges distant regions of a
knowledge network is riskier but can yield greater rewards
\cite{foster2015tradition}. These findings motivate knowledge fusion, but the
definition developed here imposes requirements beyond bibliographic novelty.

\subsection{\textbf{From Combination to Fusion}}

Let $\mathcal{K}_{1},\ldots,\mathcal{K}_{n}$ denote bodies of justified
knowledge associated with distinct domains, communities, or representational
systems. Their aggregation is the union

\begin{equation}
\begin{aligned}
\mathcal{K}_{\cup}
&=\bigcup_{i=1}^{n}\mathcal{K}_{i}.
\end{aligned}
\label{eq:knowledge_union}
\end{equation}

Aggregation increases availability but does not create a new epistemic
relation. A search engine, database, or language model may retrieve elements
from several domains and place them in one response without changing what is
known. Knowledge fusion requires an additional contribution $\kappa_F$ that is
not explicitly contained in any source domain or in their simple union:

\begin{equation}
\begin{aligned}
\mathcal{F}_{x}
(\mathcal{K}_{1},\ldots,\mathcal{K}_{n})
&\longrightarrow \kappa_F,\\
\kappa_F
&\notin \mathcal{K}_{i}
\quad \forall i,\\
\kappa_F
&\notin \mathcal{K}_{\cup}.
\end{aligned}
\label{eq:fusion_novel_relation}
\end{equation}

Here, $\mathcal{F}_{x}$ is the fusion process realized by machine $x$, and
$\kappa_F$ is a candidate relation, explanation, hypothesis, or conceptual
structure. The final condition means that $\kappa_F$ cannot be recovered as an
explicit claim already present in the aggregated sources. Its components may
be present, but their consequential connection is not.

This distinction separates five increasingly demanding operations. Retrieval
locates an existing item. Aggregation places items together. Synthesis organizes
existing claims into a coherent account. Recombination rearranges known
elements into a configuration that may be novel in form. Fusion constructs a
relation that changes what can be explained, predicted, tested, or created.
Only the last operation is a candidate pathway to knowledge invention.

\subsection{\textbf{Latent Relations Across Knowledge Domains}}

A latent relation exists when justified claims in separated knowledge domains
jointly support an implication not represented within any participating
domain. The simplest structure is the open-discovery pattern

\begin{equation}
\begin{aligned}
A &\longrightarrow B,\\
B &\longrightarrow C,\\
A &\overset{?}{\longrightarrow} C,
\end{aligned}
\label{eq:abc_fusion}
\end{equation}

where the solid relations are established and the dashed relation is an
unrecognized candidate inferred through the shared intermediary $B$. This
structure captures the foundation of classical literature-based discovery
\cite{swanson1986fish}, but knowledge fusion need not be limited to a
three-node chain. It may connect causal graphs, mathematical structures,
experimental anomalies, measurement techniques, or explanatory models across
many domains.

The importance of a fusion does not follow from distance alone. Two remote
concepts may be unrelated, and two adjacent concepts may support a profound
inference. Empirical studies of scientific production indicate that successful
high-impact work often balances conventional foundations with atypical
combinations \cite{uzzi2013atypical}. A fusion system must therefore preserve
enough established structure to support inference while exploring relations
that prevailing disciplinary organization makes unlikely to be examined.

Let $d(\mathcal{K}_{i},\mathcal{K}_{j})$ measure conceptual or bibliographic
distance, and let $J(\kappa_F)$ represent the justification available for a
candidate fusion. Then neither novelty nor distance alone is sufficient:

\begin{equation}
\begin{aligned}
d(\mathcal{K}_{i},\mathcal{K}_{j})\gg 0
&\not\Rightarrow J(\kappa_F)>0,\\
J(\kappa_F)>0
&\not\Rightarrow
\operatorname{Consequential}(\kappa_F)=1.
\end{aligned}
\label{eq:distance_not_fusion}
\end{equation}

The first relation rejects arbitrary cross-domain analogy. The second rejects
valid but trivial connections. Knowledge fusion requires both defensible
support and consequential epistemic change.

\subsection{\textbf{Fusion as the Exposure of an Unknown Unknown}}

Knowledge fusion becomes relevant to superintelligence when it reveals not
only a new answer but a previously unrecognized object of inquiry. Suppose
$\mathcal{K}_{a}$ and $\mathcal{K}_{b}$ contain individually coherent models,
yet their joint implications produce a contradiction, unexplained regularity,
or missing causal variable. A machine may use that tension to formulate a new
question $q_F$:

\begin{equation}
\begin{aligned}
\mathcal{F}_{x}(\mathcal{K}_{a},\mathcal{K}_{b})
&\longrightarrow q_F,\\
q_F
&\notin \mathcal{Q}_{H,t},\\
\operatorname{Ground}(q_F,
\mathcal{K}_{a}\cup\mathcal{K}_{b})
&=1.
\end{aligned}
\label{eq:fusion_unknown_unknown}
\end{equation}

The question is an unknown unknown before fusion because humanity has not
represented it within the recognized research agenda. Fusion changes that
state. It exposes the question, provides its initial justification, and creates
a target for discovery or theory construction. Computational work that
identifies previously unspecified state variables from observations provides a
partial example of machines constructing new representational candidates
\cite{chen2022variables}. The stronger fusion criterion additionally requires
that the machine recognize the cross-domain significance of the constructed
object.

Scientific search is shaped by human incentives and by the organization of
existing knowledge. Network analyses show that researchers preferentially
pursue established relationships even when riskier, more innovative searches
have greater potential value \cite{foster2015tradition}. Computational systems
can search differently. Models of collective discovery suggest that strategic
selection of experiments can accelerate the exploration of a scientific
knowledge network \cite{rzhetsky2015experiments}. A fusion-oriented machine
could extend this principle by selecting not only experiments but also distant
knowledge regions whose joint structure merits investigation.

\subsection{\textbf{Fusion-Driven Knowledge Invention}}

A candidate fusion becomes knowledge invention only after it is developed into
a consequential epistemic contribution and validated. Let
$\operatorname{FKI}_{x}(\kappa_F)$ denote fusion-driven knowledge invention by
machine $x$. Then

\begin{equation}
\begin{aligned}
\operatorname{FKI}_{x}(\kappa_F)
\Longleftrightarrow{}&
\operatorname{Fuse}_{x}
(\mathcal{K}_{1},\ldots,\mathcal{K}_{n})\\
&\land \kappa_F\notin\mathcal{K}_{H,t}\\
&\land
\operatorname{Originate}_{x}(\kappa_F\mid W)\\
&\land \operatorname{Consequential}(\kappa_F)\\
&\land \operatorname{Validate}(\kappa_F)=1.
\end{aligned}
\label{eq:fusion_knowledge_invention}
\end{equation}

Equation~\eqref{eq:fusion_knowledge_invention} excludes a system that merely
retrieves a previously published interdisciplinary connection. It also excludes
a system that proposes a novel connection without evidence, repeats a human
specified synthesis, or generates an interesting analogy with no explanatory
or predictive consequence. The machine must originate the essential relation
and that relation must survive independent, domain-appropriate evaluation.

Validation depends on the type of fused contribution. A causal hypothesis may
require controlled intervention; a mathematical bridge may require proof; a
unified explanatory model may require superior prediction across the source
domains; and a proposed technological principle may require successful
realization. AI-enabled science already distributes hypothesis generation,
modelling, experimental planning, and analysis across interacting systems
\cite{wang2023scientific}. Fusion-driven knowledge invention adds the
requirement that the cross-domain relation itself originate from the machine
rather than be specified by the human workflow.

\subsection{\textbf{Failure Modes and Evidentiary Safeguards}}

Knowledge fusion is especially vulnerable to outputs that appear profound
because they connect distant concepts. At least six failure modes must be
distinguished from genuine fusion. First, \emph{lexical overlap} mistakes shared
terminology for shared mechanism. Second, \emph{analogy inflation} treats a
structural resemblance as evidence of causal identity. Third,
\emph{correlation substitution} presents statistical association as
explanation. Fourth, \emph{ontology mismatch} combines concepts whose meanings
or measurement assumptions are incompatible. Fifth, \emph{hidden retrieval}
reproduces a connection present in the training corpus or literature while
presenting it as machine-originated. Sixth, \emph{novelty without consequence}
generates a defensible but scientifically unimportant relation.

For a candidate $\kappa_F$, a minimum evidentiary record should therefore
contain

\begin{equation}
\begin{aligned}
\mathcal{R}(\kappa_F)=
\{&\operatorname{Sources},
\operatorname{Provenance},
\operatorname{Derivation},\\
&\operatorname{NoveltySearch},
\operatorname{Predictions},
\operatorname{Validation}\}.
\end{aligned}
\label{eq:fusion_record}
\end{equation}

Source and provenance records identify the knowledge used. Derivation records
show how the new relation follows from or transforms those inputs. Novelty
search tests whether the relation already exists in human knowledge.
Predictions state what should be observed if the fusion is correct. Validation
records preserve the independent evidence used to accept, reject, or revise the
claim. This record does not eliminate epistemic uncertainty, but it makes the
attribution of machine-originated fusion auditable.

Knowledge fusion thus reveals what separated bodies of knowledge imply when
connected. It can expose an unknown unknown and provide the conceptual
structure from which knowledge invention develops. Knowledge discovery, the
subject of the next section, addresses a different pathway: identifying a
phenomenon, mechanism, entity, or regularity that existing knowledge does not
yet adequately represent, even after its known components have been connected.
